\section{Complex references}\label{sec:complex-references}

The activities of the renormalization group must be analytic in the background on a complex domain (Definition~\ref{def:graded}). So the step evaluates the operators of Part~IX, the interior estimates and the weakened background map at \emph{complex} references. Sections~\ref{sec:interior}--\ref{sec:weak-background} were written for real backgrounds. This section records where reality was used, and shows that every statement needed later extends to small complex perturbations. The multiscale averaged backgrounds of Part~IV need a separate argument, given in Section~\ref{sec:AB-proof} (Proposition~\ref{prop:AB}). The results of Part~IX are extended in Section~\ref{sec:complex-ext}.

\paragraph{Complex references.} Let $U_r$ be a real $(M,\eps)$-regular background. A \emph{complex reference} of size $\delta_c$ is $U=e^{\ii\eta A'}U_r$, with $A'$ a $\mathfrak g^c$-valued field such that, in the local real gauges of Definition~\ref{def:regular} for $U_r$,
\[
|A'|+|\nabla A'|+\sup_y[\nabla\tilde A'_y]_{\alpha,\Lam_1(c_y)}<\delta_c .
\]
\begin{itemize}
\item \emph{Relation to Proposition~\ref{prop:reg-complex}.} That proposition writes $U\,e^{\ii\eta\,\delta\!A}$ and uses ordinary derivatives. The two forms are equivalent, with $\delta\!A=R(U_r)^{-1}A'$. The two sizes are comparable: there is $\eps_{\rm cc}>0$, depending only on $d$, $\alpha$ and $G$, such that for $\eps\le\eps_{\rm cc}$ the $\delta$ of Proposition~\ref{prop:reg-complex} and the size of $A'$ above bound each other up to the factor $C_{\rm cmp}=4$. Indeed, in the gauge of Definition~\ref{def:regular} on a cube $\hat\Lam$, where $U_r=e^{\ii\eta A_r}$, the two fields differ by the factor $R(e^{\ii\eta A_r})$, which is within $4\eta\eps$ of $1$, with difference quotients and H\"older quotients of size $C\eta\eps$. Covariant and flat gradients differ by $C\eps|A'|$, by \eqref{eq:cov-vs-flat}, and the change between the gauges of the blocks and of the cubes costs at most $C\eps$ times the $C^1$ norm, as in the proof of Theorem~\ref{thm:ID}(b). Finally, two points at $\ell^\infty$-distance less than $1$ lie in a common cube $\Lam_1(c_y)$, and pairs at distance at least $1$ have H\"older quotients at most twice the sup of the gradient.
\item \emph{The local gauge.} In a local gauge of $U_r$ we have $U=e^{\ii\eta A_0}$, with $A_0=A'\oplus A_r$ combined bondwise by the Baker--Campbell--Hausdorff formula. Then $A_0$ is complex and satisfies (B) with parameter $C(\eps+\delta_c)$.
\item \emph{Potentials.} Local potentials at a complex reference are taken in $\cW^{\rm an}_{\rm loc}$ (Section~\ref{sec:reg-gauge}), and satisfy the bounds of $\cW_{\rm loc}(\eps+C\delta_c)$ in the real gauges of $U_r$. The non-local potential $W_2^X$ satisfies the hypotheses of Proposition~\ref{prop:reg-complex} for existence, $\Rea$-coercivity and decay; its comparison of domains is Proposition~\ref{prop:dirichlet-hessian}(d). Only its $L^2$ statements are used (Section~\ref{sec:localization}). In Section~\ref{sec:weak-background}, $W=0$.
\item \emph{Graded domains.} A point of a graded domain $\mathcal U(P;\theta_r,\delta_1,\delta_0)$ fixes the complex part only near $P$. It is extended by $U_r$ outside a neighbourhood, with a smooth cutoff at unit scale (Remark~\ref{rem:bg-to-graded}). This gives a complex reference with $\delta_c\le3C_{\rm ext}(\delta_0+\delta_1)$; for a point of the domain of step $k+1$, read in the units of step $k$, $\delta_c\le2C_{\rm ext}(L^{-1}\delta_0+L^{-2}\delta_1)$ (Remark~\ref{rem:domain-constants}(ii)). The real part has only curvature and current below $\theta_r$. Its $(M,B_5\theta_r)$-regularity comes from Corollary~\ref{cor:regularity}, with $B_5$ depending on $L$ and $M$. Its current is at most $\eps$ if $\theta_r\le\eps$; Section~\ref{sec:step-quantities} adds $|J_{U_r}|\le\eps$ to the definition of a complex reference, for the coercivity of the Hessian.
\end{itemize}

\subsection{Averaged backgrounds at complex references}\label{sec:AB-proof}

\begin{readingbox}[Recall from Part~IV]
This subsection uses the multiscale averaging of Sections~\ref{sec:averaging} and~\ref{sec:adapted}. In the units of step $k$, the lattice $\T^{(j)}$ has spacing $\delta_j=L^{j-k}$, with $\T^{(0)}=\Tf$ and $\T^{(k)}=\T$ (Section~\ref{ss:21.1}). Here $\delta_j$ is a spacing, unrelated to the domain constants $\delta_0,\delta_1$ of Remark~\ref{rem:domain-constants}.
\begin{itemize}
\item \emph{Group average.} $\ell$ is the logarithm chart of $G$, of radius $c_G$, and $\kappa(u)=\|u\|\|u^{-1}\|$ for $u\in G^c$ (Section~\ref{ss:21.2}). For a family $\mathbf U$ in $\cD_{r_M}$, where $\cD_r=\{\operatorname{diam}\mathbf U<r\}$, the group average $\Av(\mathbf U)$ is the solution of \eqref{eq:group-average} given by Lemma~\ref{lem:group-average}(a); $r_M$, $C_M$ and the weights $\omega_\Gamma$ of its derivative are those of that lemma.
\item \emph{Ba\l aban's averaging.} $Y(y,x)$ is the averaged contour variable, the group average over the contour family $G(y,x)$, and $\Lambda_{c,x}(Y)$ are the averaged loops (Definition~\ref{def:avg-contour}); their gauge covariance is \eqref{eq:avg-contour-gauge}. The one-step average $\mathfrak a_j$ is \eqref{eq:one-step}, $\mathfrak a^{(j)}=\mathfrak a_j\circ\mathfrak a^{(j-1)}$ and $U^{(j)}=\mathfrak a^{(j)}(U)$ (Definition~\ref{def:one-step}). $\mathfrak U_j(a)$ is the set of real configurations on $\T^{(j)}$ with curvature at most $a$ (Definition~\ref{def:curvature}), and $U^{(j)}\in\mathfrak U_j(2a)$ for $U\in\mathfrak U_0(a)$ with $a\le c_1$ (Proposition~\ref{prop:avg-regular}).
\item \emph{Gauge-adapted averaging.} A relative configuration $X$ at scale $j$ stands for $XW_j$ (Section~\ref{ss:22.1}). $\tau_X$ is the block average of its averaged relative contour variables, $\tilde{\mathfrak a}_j(X)=\overline X^{(\tau_X^{-1})}$ the modified one-step average, $\mathcal Q_j=\tilde{\mathfrak a}_j\circ\mathcal Q_{j-1}$ the gauge-adapted averaging, and $\mathsf w_j$ the correspondence transformation, built with the extensions $\mathcal E^{(i)}$ (Definition~\ref{def:adapted}).
\item \emph{Linearizations.} $\tilde{\mathfrak a}_j(e^{\ii\delta_{j-1}B})=e^{\ii\delta_j(l_jB+c_j(B))}$ and $\tau_X=e^{\ii\delta_j(\gamma_jB+c^\tau_j(B))}$, with $l_j=q_j+r_j$ and the constant $c_2$ (Lemma~\ref{lem:one-step-expansion}); the averaged relative transports $\lambda_W(y,x)$ (Definition~\ref{def:avg-transport}); $\bQ_j=l_j\cdots l_1$, the nonlinear part $\Cav_j$ and the constant $c_3$ (Proposition~\ref{prop:Q-expansion}); $q'_j$ and $Q'_k$ (Section~\ref{ss:22.6}, Proposition~\ref{prop:linearized-averages}).
\end{itemize}
Below, $h_j=\tau_0\cdots\tau_{j-1}$ is a \emph{complex} gauge transformation on $\T^{(j)}$, built from the relative configuration of a complex reference (display after the Setting). It is not the partition of unity $h_j$ of Sections~\ref{sec:weakenable}--\ref{sec:weak-background}, nor the threshold $h_k=p_kg_k$ of Part~X.
\end{readingbox}

Proposition~\ref{prop:Q-expansion} and Proposition~\ref{prop:linearized-averages} assume a real $U_0$, with real averaged backgrounds $W_j\in\mathfrak U_j(2a)$ (Proposition~\ref{prop:avg-regular}). Lemma~\ref{lem:one-step-expansion} uses this through:
\begin{itemize}
\item the unitarity of the transports $R(W(\gamma))$ along contours and $R(W(y,x))$ along averaged contour variables;
\item the families of the background in $\cD_{6a\delta_j^2}$ (Lemma~\ref{lem:stokes-avg}(i)), and through them the weights of the averaged relative transports, with $\sum_\Gamma|\omega_\Gamma|\le1+Ca\delta_j^2$ (Lemma~\ref{lem:avg-transport}(c));
\item the bound on the averaged loops $h_x$ and on $\Eloop$ (Lemma~\ref{lem:one-step-basic}(a));
\item Lemma~\ref{lem:avg-local-gauge}.
\end{itemize}
The underlying lemmas are proved only for real configurations: Lemma~\ref{lem:log-estimates}(b), with the bi-invariant metric on $G$; Lemmas~\ref{lem:stokes} and~\ref{lem:stokes-avg}(i)--(iii); and Lemma~\ref{lem:one-step-regularity}, which uses a real gauge and gauge invariance of $|F|$. These results feed the bounds $|\bQ A(c)|\le C\|A\|$ and $|Q'\lambda|\le C\|\lambda\|$, and the bounded transposes, used in Theorem~\ref{thm:ID} and in Part~IX. A naive per-scale bound on complex transports would compound to $2^{j}$ over the scales.

We therefore work relative to the real reference $U_r$. Two observations make this work.
\begin{itemize}
\item The averaged backgrounds at $U$ are gauge transforms of $X_jW_j(U_r)$, with $X_j=\mathcal Q_j(U')$, by explicit complex gauge transformations $h_j$ (Lemma~\ref{lem:AB-rep}). This is proved directly from the one-step definitions, with the complex equivariance of the group average (Lemma~\ref{lem:group-average}(b)). The axial gauge, Lemma~\ref{lem:axial}(b),(c) and Theorem~\ref{thm:correspondence}(a),(b),(e),(f) are not used at complex configurations.
\item At the points of $\T^{(j)}$, $h_j$ has size $C\delta_j|A'|$, not $C|A'|$. It is a product of the variables $\tau_{X_i}$, $i<j$. Each $\tau_{X_i}$ is the exponential of an average of logarithms of averaged relative contour variables. These are group averages of relative holonomies along the contours of the families $G(y,x)$, of length at most $2(L-1)\delta_i<2\delta_{i+1}$, and by the Lipschitz bound of Lemma~\ref{lem:group-average}(e) they are of the same size. Moreover $\sum_{i<j}\delta_{i+1}\le2\delta_j$. The scale-$j$ gradient of $h_j$ carries the factor $\delta_j^{-1}$, so it is bounded by the sup norm. Hence the relative fields are $O(|A'|)$ uniformly in $j$, and the losses per scale are summable.
\end{itemize}
Only the sup norm of $A'$ is used, and the smallness of $\delta_c$ depends on $d$, $L$ and $G$, but not on $M$.

\paragraph{Setting.} Let $U_r$ be real and $(M,\eps)$-regular with $\eps\le\eps_0$. By the proof of Corollary~\ref{cor:parts-apply}, $U_r\in\mathfrak U_0(a)$ with $a\le C\eps$, where $C$ does not depend on $M$ (Lemma~\ref{lem:curvature}). The proofs of Sections~\ref{sec:averaging} and~\ref{sec:adapted} do not use $M$. So $c_3$ and the constants $C,c$ of this subsection depend only on $d$, $L$ and $G$. Let $c'_3\le c_3$ be such a constant, decreased below where needed, and assume $a\le c'_3$. Put $W^r_j=W_j(U_r)=U_r^{(j)}$. These are real and lie in $\mathfrak U_j(2a)$ by Proposition~\ref{prop:avg-regular}.

Let $U=U'U_r$ (bondwise), with $U'=e^{\ii\eta A'}$ and $A'$ a $\mathfrak g^c$-valued field. Under a real gauge transformation $u$, $U_r\mapsto U_r^u$ and $A'(b)\mapsto R(u(b_-))A'(b)$. So $|A'|_\infty$ is gauge invariant, and the condition $|A'|<\delta_c$ in the local gauges of $U_r$ is the condition $|A'|_\infty<\delta_c$. The derivative bounds on $A'$ in the definition of complex references are not used in this subsection. Let $W_j(U)=U^{(j)}=\mathfrak a^{(j)}(U)$, whenever defined.

We use the gauge-adapted averaging of Section~\ref{sec:adapted} relative to the \emph{real} reference $U_r$, so there $W_j=W^r_j$. Applied to the relative configuration $U'$ it gives
\[
X_j=\mathcal Q_j(U')=e^{\ii\delta_jB_j},\qquad \tau_i=\tau_{X_i}\ \ (\text{on }\T^{(i+1)}),\qquad h_j(y)=\tau_0(y)\cdots\tau_{j-1}(y)\ \ (y\in\T^{(j)}),
\]
with $h_0=1$. By Theorem~\ref{thm:correspondence}(c), which holds for complex relative configurations (remark after that theorem), $h_j=\mathsf w_j(U')^{-1}$ at the points of $\T^{(j)}$. Recall that $\tilde{\mathfrak a}_j(X)=\overline X^{(\tau_X^{-1})}$ (Definition~\ref{def:adapted}); on the domain of Lemma~\ref{lem:one-step-expansion} this is $\overline{X^{(g_X)}}$, by Lemma~\ref{lem:one-step-basic}(b).

Norms of linear maps of $\mathfrak g^c$, and of matrices acting on $\C^{\mathsf N}$, are operator norms, written $\|\cdot\|$. The constants
\[
C_{\rm AB}=14d+3,\qquad\kappa_{\rm AB}=16C_MC_{\rm AB},
\]
with $C_M\ge1$ of Lemma~\ref{lem:group-average}, and $C_{\rm loop}$ of Lemma~\ref{lem:AB-bounds}(d) depend only on $d$ and $G$. The contours used in the step from scale $j-1$ to scale $j$, namely those of the families $G(y,x)$, $y\in\T^{(j)}$, the segments $[x,x(c)]$ and $[c]$, and their initial pieces, have at most $2L$ bonds of $\T^{(j-1)}$ (Section~\ref{sec:averaging}). Recall $\delta_j=L^{j-k}$, so that
\begin{equation}\label{eq:AB-scales}
\sum_{i\le k}\delta_i\le2,\qquad\sum_{i<j}\delta_{i+1}=\delta_j\sum_{m<j}L^{-m}\le2\delta_j .
\end{equation}

\begin{lemma}[Covariance under complex gauge transformations]\label{lem:AB-cov}
Let $j\ge1$ and let $Y$ be a configuration on $\T^{(j-1)}$ such that all families $\{Y(\Gamma)\}_{\Gamma\in G(y,x)}$, $y\in\T^{(j)}$, $x\in B_j(y)$, lie in $\cD_{r_M/2}$, and $|\Lambda_{c,x}(Y)-1|\le\frac12c_G$ for all bonds $c$ of $\T^{(j)}$, of both orientations, and all $x\in B_j(c_-)$. Let $g:\T^{(j-1)}\to G^c$ with $\kappa(g(y))\le2$ for $y\in\T^{(j)}$. Then $\mathfrak a_j(Y)$ and $\mathfrak a_j(Y^g)$ are defined, and $\mathfrak a_j(Y^g)=\mathfrak a_j(Y)^{g|_{\T^{(j)}}}$.

Consequently, let $X$ be a relative configuration at scale $j-1$ such that $XW_{j-1}$ satisfies the hypothesis on $Y$, for some background $W_{j-1}$ with $W_j=\mathfrak a_j(W_{j-1})$. Then $\overline{X^{(g)}}=\overline X^{(g|_{\T^{(j)}})}$.
\end{lemma}

\begin{proof}\hasdetail{lem:AB-cov}
This is the complex case of Lemma~\ref{lem:one-step-basic}(b). In detail: by Lemma~\ref{lem:group-average}(b), the families of $Y^g$ lie in $\cD_{r_M}$, and \eqref{eq:avg-contour-gauge} holds. Since $c_\pm\in\T^{(j)}$, this gives $\Lambda_{c,x}(Y^g)=g(c_-)\Lambda_{c,x}(Y)g(c_-)^{-1}$, and $Y^g([c])=g(c_-)Y([c])g(c_+)^{-1}$. Moreover $|ghg^{-1}-1|\le\kappa(g)|h-1|\le c_G$ for $h=\Lambda_{c,x}(Y)$ and $g=g(c_-)$, so Lemma~\ref{lem:log-estimates}(d) gives $\ell(ghg^{-1})=g\,\ell(h)\,g^{-1}$, and $\exp(\ii gZg^{-1})=g\,e^{\ii Z}g^{-1}$. Insert this into \eqref{eq:one-step}. For relative configurations use $X^{(g)}W_{j-1}=(XW_{j-1})^g$ and $\overline X=\mathfrak a_j(XW_{j-1})W_j^{-1}$.
\end{proof}

This is the form of Lemma~\ref{lem:one-step-basic}(b) used below. For real $g$ only the definedness of $\mathfrak a_j(Y)$ is needed.

\begin{lemma}[Representation of the averaged backgrounds]\label{lem:AB-rep}
There is $\delta_{AB}>0$, depending on $d,L,G$, such that for $a\le c'_3$ and $|A'|_\infty\le\delta_{AB}$:
\begin{enumerate}[label=(\alph*)]
\item $|\ell(\tau_i)|\le3d\,\delta_{i+1}|A'|_\infty$ on $\T^{(i+1)}$ for $i<k$. The function $\lambda_j=\ell(h_j)$ on $\T^{(j)}$ is defined and satisfies $|\lambda_j|\le7d\,\delta_j|A'|_\infty$ for $j\le k$.
\item For every $j\le k$, $W_j(U)$ is defined and
\begin{equation}\label{eq:AB-rep}
W_j(U)=\big(X_jW^r_j\big)^{h_j},\qquad\text{i.e.}\quad W_j(U)(b)=h_j(b_-)\,X_j(b)\,W^r_j(b)\,h_j(b_+)^{-1}.
\end{equation}
Equivalently, $U'^{[j]}=\mathcal Q_j(U')^{(h_j)}=\mathcal Q_j(U')^{(\mathsf w_j(U')^{-1})}$.
\end{enumerate}
\end{lemma}

\begin{proof}\hasdetail{lem:AB-rep}
(a) Proposition~\ref{prop:Q-expansion}(a) applies to the real reference $U_r$ and the complex field $A'$, and gives $|B_i|_\infty\le2|A'|_\infty$. By Lemma~\ref{lem:one-step-expansion}, $\ell(\tau_i)=\delta_{i+1}\big(\gamma_{i+1}B_i+c^\tau_{i+1}(B_i)\big)$ with $|\gamma_{i+1}B_i|\le2d|A'|$ and $|c^\tau_{i+1}(B_i)|\le4C\delta_{i+1}|A'|^2$. This gives the first bound for $|A'|\le d/(4C)$. By \eqref{eq:AB-scales}, $s=\sum_{i<j}|\ell(\tau_i)|\le6d\delta_j|A'|$. Lemma~\ref{lem:log-estimates}(c) gives $|\lambda_j|\le s+Cs^2\le7d\delta_j|A'|$ for $|A'|$ small.

(b) We use induction on $j$. For $j=0$ we have $W_0(U)=U=U'U_r=X_0W^r_0$ and $h_0=1$. Let $j\ge1$, assume \eqref{eq:AB-rep} at $j-1$, and put $Y=X_{j-1}W^r_{j-1}$. By Proposition~\ref{prop:Q-expansion}(a), $X_{j-1}=e^{\ii\delta_{j-1}B_{j-1}}$ with $|B_{j-1}|_\infty\le2|A'|_\infty\le c_2$, after decreasing $\delta_{AB}$. So Lemma~\ref{lem:one-step-expansion}, at the real background $W^r_{j-1}\in\mathfrak U_{j-1}(2a)$ and with the complex Lie field $B_{j-1}$, shows that $X_{j-1}$ lies in the averaging domain at scale $j$: the families of $Y$ lie in $\cD_{r_M/2}$, and its averaged loops satisfy $|\Lambda_{c,x}(Y)-1|\le\frac12c_G$. These loops are built from group averages of families; the proof of that lemma bounds them through the Lipschitz bound of Lemma~\ref{lem:group-average}(e), and Lemma~\ref{lem:stokes-avg}(iv) gives a bound of the same kind. In particular $\mathfrak a_j(Y)$ is defined. By (a), $h_{j-1}=e^{\ii\lambda_{j-1}}$ with $\kappa(h_{j-1})\le e^{2|\lambda_{j-1}|}\le2$. By the induction hypothesis $W_{j-1}(U)=Y^{h_{j-1}}$, so Lemma~\ref{lem:AB-cov} shows that $W_j(U)=\mathfrak a_j(W_{j-1}(U))$ is defined and
\[
W_j(U)=\mathfrak a_j(Y)^{h_{j-1}|_{\T^{(j)}}} .
\]
By the definitions of the relative average and of $\tilde{\mathfrak a}_j$,
\[
\mathfrak a_j(Y)=\overline{X_{j-1}}\,W^r_j,\qquad \overline{X_{j-1}}=\tilde{\mathfrak a}_j(X_{j-1})^{(\tau_{j-1})}=X_j^{(\tau_{j-1})},
\]
i.e.\ $\mathfrak a_j(Y)=(X_jW^r_j)^{\tau_{j-1}}$. Since $(Z^g)^{g'}=Z^{g'g}$ and $h_{j-1}\tau_{j-1}=h_j$ on $\T^{(j)}$, \eqref{eq:AB-rep} follows at $j$. The last statement is \eqref{eq:AB-rep} rewritten with the relative gauge action, together with $h_j=\mathsf w_j(U')^{-1}$ on $\T^{(j)}$.
\end{proof}

\begin{lemma}[Averaged backgrounds at complex references]\label{lem:AB-bounds}
Possibly after decreasing $\delta_{AB}$ and $c'_3$ (depending on $d,L,G$), the following holds for $a\le c'_3$, $|A'|_\infty\le\delta_c\le\delta_{AB}$ and $0\le j\le k$.
\begin{enumerate}[label=(\alph*)]
\item $W_j(U)=e^{\ii\delta_j\tilde B_j}\,W^r_j$ bondwise, where
\begin{equation}\label{eq:AB-field}
\tilde B_j=B_j-D_{W^r_j}\lambda_j+\beta_j,\qquad|\beta_j|_\infty\le C\delta_j|A'|^2_\infty,\qquad|\tilde B_j|_\infty\le C_{\rm AB}|A'|_\infty .
\end{equation}
Moreover $\big|\tilde B_j-\bQ_j(U_r)A'+D_{W^r_j}\boldsymbol\Gamma_jA'\big|_\infty\le C|A'|_\infty^2$, where $\boldsymbol\Gamma_jA'=\sum_{i<j}\delta_{i+1}\gamma_{i+1}\bQ_i(U_r)A'$ is evaluated at the points of $\T^{(j)}$.
\item $\tilde B_j$ is holomorphic in $A'$ on $\{|A'|_\infty<\delta_{AB}\}$, and real for real $A'$. Under a real gauge transformation $u$ of $\Tf$ ($U_r\mapsto U_r^u$, $A'(b)\mapsto R(u(b_-))A'(b)$) we have $\tilde B_j(b)\mapsto R(u(b_-))\tilde B_j(b)$. In particular (a) holds in every real gauge, e.g.\ in the local gauges of $U_r$.
\item Let $\gamma$ be a contour of $n$ bonds of $\T^{(j)}$ with $s=C_{\rm AB}n\delta_j|A'|_\infty\le s_0$, where $s_0$ depends only on $\mathsf N$. Then $W_j(U)(\gamma)=e^{\ii Z_\gamma}W^r_j(\gamma)$ with $|Z_\gamma|\le2s$. (Here $Z_\gamma$ is the Baker--Campbell--Hausdorff series of $\mathfrak g^c$-elements of Lemma~\ref{lem:log-estimates}(c), which lies in $\mathfrak g^c$; so no radius $c_G$ enters.) Hence
\[
\|R(W_j(U)(\gamma))^{\pm1}\|\le e^{4s},\qquad\|R(W_j(U)(\gamma))-R(W^r_j(\gamma))\|\le e^{4s}-1 .
\]
Moreover, for $j\ge1$, $y\in\T^{(j)}$ and $x\in B_j(y)$, the averaged contour variable $W_{j-1}(U)(y,x)$ is defined, and
\[
W_{j-1}(U)(y,x)=e^{\ii Z_{y,x}}\,W^r_{j-1}(y,x),\qquad|Z_{y,x}|\le4C_MC_{\rm AB}\delta_j|A'|_\infty .
\]
Applied at index $j-1$ to the contours of at most $2L$ bonds of $\T^{(j-1)}$ used in the step $j-1\to j$, where $4s\le8C_{\rm AB}\delta_j\delta_c$, and to the averaged contour variables, where $2|Z_{y,x}|\le8C_MC_{\rm AB}\delta_j\delta_c$, this gives these bounds for both kinds of transports with $e^{4s}$ replaced by $e^{\kappa_{\rm AB}\delta_j\delta_c/2}$. Products of one transport of each kind have norm at most $e^{\kappa_{\rm AB}\delta_j\delta_c}$.
\item For $j\ge1$, a bond $c$ of $\T^{(j)}$ and $x\in B_j(c_-)$, let
\[
H_x(U)=\ell\big(\Lambda_{c,x}(W_{j-1}(U))\big),\qquad \Eloop(U)=L^{-d}\sum_xH_x(U),
\]
the logarithms of the averaged loops of \eqref{eq:one-step} and their average, so that $W_j(U)(c)=e^{\ii \Eloop(U)}W_{j-1}(U)([c])$. Then
\begin{gather*}
|H_x(U)-H_x(U_r)|,\ \ |\Eloop(U)-\Eloop(U_r)|\le C_{\rm loop}\delta_j|A'|_\infty,\\
|H_x(U)|\le16a\delta_j^2+C_{\rm loop}\delta_j|A'|_\infty,\qquad|\Eloop(U)|\le16a\delta_j^2+C_{\rm loop}\delta_j|A'|_\infty .
\end{gather*}
\end{enumerate}
\end{lemma}

\begin{proof}\hasdetail{lem:AB-bounds}
(a) By Lemma~\ref{lem:AB-rep}, $V_j=W_j(U)(W^r_j)^{-1}=X_j^{(h_j)}$, i.e.
\[
V_j(b)=e^{\ii\lambda_j(b_-)}\,e^{\ii\delta_jB_j(b)}\,e^{-\ii R(W^r_j(b))\lambda_j(b_+)} .
\]
The three exponents have total size $s\le(14d+2)\delta_j|A'|\le\frac14$, by Lemma~\ref{lem:AB-rep}(a) and Proposition~\ref{prop:Q-expansion}(a); here we used that $R(W^r_j)$ is unitary. By Lemma~\ref{lem:log-estimates}(c), $\ell(V_j(b))=\lambda_j(b_-)+\delta_jB_j(b)-R(W^r_j(b))\lambda_j(b_+)+\delta_j\beta_j(b)$ with $\delta_j|\beta_j|\le Cs^2$. This is \eqref{eq:AB-field}, with $D_{W^r_j}\lambda(b)=\delta_j^{-1}\big(R(W^r_j(b))\lambda(b_+)-\lambda(b_-)\big)$ as in Section~\ref{sec:adapted}. By Lemma~\ref{lem:AB-rep}(a), $|D_{W^r_j}\lambda_j|\le\delta_j^{-1}\big(|\lambda_j(b_+)|+|\lambda_j(b_-)|\big)\le14d|A'|$. Hence $|\tilde B_j|\le(2+14d+C\delta_j|A'|)|A'|\le C_{\rm AB}|A'|$.

For the linear part, use $B_j=\bQ_jA'+\Cav_j(A')$ with $|\Cav_j(A')|\le C|A'|^2$ (Proposition~\ref{prop:Q-expansion}(b)). Further, $\lambda_j=\sum_{i<j}\ell(\tau_i)+O(\delta_j^2|A'|^2)$ and $\ell(\tau_i)=\delta_{i+1}\gamma_{i+1}\bQ_iA'+O(\delta_{i+1}|A'|^2)$. By \eqref{eq:AB-scales}, $|\lambda_j-\boldsymbol\Gamma_jA'|\le C\delta_j|A'|^2$, and $D_{W^r_j}$ costs at most $2\delta_j^{-1}$.

(b) $B_j$ and $\tau_i$ are holomorphic in $A'$ by Proposition~\ref{prop:Q-expansion} and Lemma~\ref{lem:one-step-expansion}. Then $\lambda_j$ and $\tilde B_j$ are holomorphic, since they are obtained from these by $\ell$, products and Lemma~\ref{lem:log-estimates}(c). For real $A'$ all configurations involved are real. For covariance, $U^u=e^{\ii\eta R(u_-)A'}U_r^u$. Lemma~\ref{lem:one-step-basic}(b) with real $g$, applied at each scale, where the averages are defined by Lemma~\ref{lem:AB-rep}(b), gives $W_j(U^u)=W_j(U)^{u|_{\T^{(j)}}}$, and likewise $W_j(U_r^u)=(W^r_j)^{u|_{\T^{(j)}}}$. Hence $V_j(b)\mapsto u(b_-)V_j(b)u(b_-)^{-1}$.

(c) $W_j(U)(\gamma)=(V_j)_{W^r_j}(\gamma)\,W^r_j(\gamma)$, and $(V_j)_{W^r_j}(\gamma)=\prod_iR(W^r_j(\gamma_{<i}))\big(V_j(b_i)\big)$. For negatively oriented bonds, $V_j(-b)=R(W^r_j(b))^{-1}(V_j(b)^{-1})$. Since $R(W^r_j)$ is unitary, every factor is $e^{\ii\delta_jZ_i}$ with $|Z_i|\le C_{\rm AB}|A'|$. Lemma~\ref{lem:log-estimates}(c) gives $|Z_\gamma|\le s+Cs^2\le2s$ for $s\le s_0=\min(\frac14,C^{-1})$. By \eqref{eq:Ad-bound}, $\|R(e^{\pm\ii Z})\|\le e^{2|Z|}$ and $\|R(e^{\ii Z})-1\|\le e^{2|Z|}-1$, and $R(W^r_j(\gamma))$ is unitary.

For the averaged contour variables, the contours $\Gamma\in G(y,x)$ have at most $2(L-1)$ bonds of $\T^{(j-1)}$, with $\delta_{j-1}=\delta_j/L$. By the first part at index $j-1$, $W_{j-1}(U)(\Gamma)=e^{\ii Z_\Gamma}W^r_{j-1}(\Gamma)$ with $|Z_\Gamma|\le4C_{\rm AB}\delta_j|A'|$. The real family $\{W^r_{j-1}(\Gamma)\}$ lies in $\cD_{6a\delta_j^2}\subset\cD_{r_M/2}$ (Section~\ref{sec:adapted}, setting), and $|Z_\Gamma|\le r_M/16$ after decreasing $\delta_{AB}$. By Lemma~\ref{lem:group-average}(a),(e) and the proof of (e), the family $\{W_{j-1}(U)(\Gamma)\}$ lies in $\cD_{r_M}$, its group average is the averaged contour variable $W_{j-1}(U)(y,x)$, and $|\ell(W_{j-1}(U)(y,x)W^r_{j-1}(y,x)^{-1})|\le C_M\max_\Gamma|Z_\Gamma|$. The bounds on the transports follow from \eqref{eq:Ad-bound}, since $R(W^r_{j-1})$ is unitary and $8C_{\rm AB}\le\frac12\kappa_{\rm AB}$. They need $2C_{\rm AB}\delta_j\delta_c\le s_0$, which holds after decreasing $\delta_{AB}$.

(d) The averaged loop $\Lambda_{c,x}(W_{j-1}(U))$ is a product of four factors: the averaged contour variables $W_{j-1}(U)(c_-,x)$ and $W_{j-1}(U)(x(c),c_+)=W_{j-1}(U)(c_+,x(c))^{-1}$, and the transports along $[x,x(c)]$ and along $[c]^{-1}$, two straight contours of $L$ bonds of $\T^{(j-1)}$. By (c) at index $j-1$, each factor is $e^{\ii\zeta}$ times the corresponding factor at $U_r$, or that factor times $e^{\ii\zeta}$, with $|\zeta|\le4C_MC_{\rm AB}\delta_j|A'|$ for the averaged contour variables and $|\zeta|\le2C_{\rm AB}\delta_j|A'|$ for the straight contours. The factors at $U_r$ are real. Moving the factors $e^{\ii\zeta}$ to the left conjugates them by real elements, which preserves their norms. So
\[
\Lambda_{c,x}(W_{j-1}(U))=e^{\ii\eta_1}e^{\ii\eta_2}e^{\ii\eta_3}e^{\ii\eta_4}\,e^{\ii H},\qquad\sum_i|\eta_i|\le12C_MC_{\rm AB}\delta_j|A'|,
\]
with $H=H_x(U_r)$. By Lemma~\ref{lem:one-step-basic}(a) and Proposition~\ref{prop:avg-regular}, $|H|\le16a\delta_j^2\le16c'_3$. Let $\varphi(t)=\ell(e^{\ii t\eta_1}\cdots e^{\ii t\eta_4}e^{\ii H})-H$. For $|t|\le r_0/\sum_i|\eta_i|$, with $r_0$ small depending on $\mathsf N$ (the logarithm is again a Baker--Campbell--Hausdorff series of $\mathfrak g^c$-elements, so $c_G$ does not enter), it is holomorphic and bounded by a constant $K_{\rm S}(\mathsf N)$, and $\varphi(0)=0$; this uses $16c'_3\le r_0$. The Schwarz lemma gives $|\varphi(1)|\le(K_{\rm S}/r_0)\sum_i|\eta_i|$. This is the bound with $C_{\rm loop}=12C_MC_{\rm AB}K_{\rm S}/r_0$. Average over $x$ for $\Eloop$.
\end{proof}

\begin{proposition}[Averaged backgrounds at complex references, (AB)]\label{prop:AB}
There are $\delta_{AB}>0$ and $c'_3\le c_3$, depending only on $d,L,G$ (through the constants of Part~IV), such that the following holds for $a\le c'_3$, every real $U_r\in\mathfrak U_0(a)$, and every complex reference $U=e^{\ii\eta A'}U_r$ with $|A'|_\infty\le\delta_c<\delta_{AB}$. All bounds are uniform in $k$, in $\eta$ and in $j\le k$. Below, relative configurations, $\tau_X$, $\tilde{\mathfrak a}_j$, $\mathcal Q_j$, $l_j$, $q_j$, $r_j$, $\gamma_j$, $q'_j$, $\bQ_k$, $Q'_k$ at the reference $U$ are defined as in Section~\ref{sec:adapted}, with the complex backgrounds $W_j(U)$ in place of $W_j$.
\begin{enumerate}[label=(\roman*)]
\item \emph{(Backgrounds.)} $W_j(U)$ is defined for $j\le k$, and Lemmas~\ref{lem:AB-rep} and~\ref{lem:AB-bounds} hold. In particular, relative to $W_j(U_r)$, $W_j(U)$ is a complex perturbation of size at most $C_{\rm AB}\delta_c$, in every real gauge.
\item \emph{(One step.)} Lemma~\ref{lem:one-step-expansion} holds at $W=W_{j-1}(U)$, with $c_2$ replaced by $\frac12c_2$. Parts (a), (c), (d) and (f) hold as stated, except the reality statement in (f); (e) is dropped, and (b) is replaced by
\begin{gather*}
|(q_jB)(c)|^2\le e^{2\kappa_{\rm AB}\delta_j\delta_c}L^{-d}\sum_xL^{-1}\sum_{b\in[x,x(c)]}|B(b)|^2,\\
|(r_jB)(c)|\le C(a\delta_j^2+\delta_j\delta_c)\max|B|,\qquad|\gamma_jB|\le de^{\kappa_{\rm AB}\delta_j\delta_c}\max|B| .
\end{gather*}
\item \emph{(Gauge-adapted averaging.)} Proposition~\ref{prop:Q-expansion} holds at the reference $U$, for complex $A$ with $|A|_\infty\le c'_3$. Parts (a), (b) and (d) hold in the same form, and ``real for real $A$'' is dropped. In (c) the kernel bounds become
\begin{gather*}
|K(c,b)|\le(1+Ca)\,e^{C_L\delta_c},\qquad|K_A(c,b)|\le Ce^{C_L\delta_c}|A|_\infty,\\
|d^2\Cav_k(A)[h,h](c)|\le Ce^{C_L\delta_c}\|h\|^2_{S_c},
\end{gather*}
and the consequences of (c) (pointwise bounds by $\|h\|_{S_c}$, sup-norm bounds of the transposes) hold with their constants multiplied by $e^{C_L\delta_c}$.
\item \emph{(Linearized averages.)} $Q'_k(U)Q'_k(U)^T=I$, with the transpose taken with respect to $\mathrm b$. Proposition~\ref{prop:linearized-averages}(d), and (f) for real $u$, hold as stated. The transports along the chains of averaged contour variables satisfy $\|T_U(y,x)\|\le e^{\kappa_{\rm AB}\delta_c}$ and $\|T_U(y,x)-T_{U_r}(y,x)\|\le e^{\kappa_{\rm AB}\delta_c}-1$. Hence $|Q'_k(U)\lambda(y)|\le e^{\kappa_{\rm AB}\delta_c}\|\lambda\|_{B(y)}$, and $Q'_k(U)^T$ is bounded on sup norms by $e^{\kappa_{\rm AB}\delta_c}$. Proposition~\ref{prop:linearized-averages}(e) holds with $\eps$ replaced by $\eps+\delta_c$. More precisely, suppose that on a cube containing $B(c_-)\cup B(c_+)$ we have $U_r=e^{\ii\eta A_r}$ in a real gauge, with $|A_r|\le\eps$ and $a\le\eps\le c'_3$. Then, in this gauge, $|K_{\bQ_k(U)}(c,b)-K_{\bQ}(c,b)|+|K_{Q'_k(U)}(y,x)-K_{Q_0}(y,x)|\le C_L(\eps+\delta_c)$, and the same holds for the transposes.
\item \emph{(Closeness to the real reference.)} With the supports of Proposition~\ref{prop:linearized-averages}(d),
\begin{gather*}
|K_{\bQ_k(U)}(c,b)-K_{\bQ_k(U_r)}(c,b)|\le C_L\delta_c,\qquad |K_{Q'_k(U)}(y,x)-K_{Q'_k(U_r)}(y,x)|\le e^{\kappa_{\rm AB}\delta_c}-1,\\
\|R(W_k(U)(b))-R(W_k(U_r)(b))\|\le e^{2C_{\rm AB}\delta_c}-1 .
\end{gather*}
So $\bQ_k(U)-\bQ_k(U_r)$, $Q'_k(U)-Q'_k(U_r)$, $D_{W_k(U)}-D_{W_k(U_r)}$ and their transposes are local operators of size $C_L\delta_c$.
\item \emph{(Analyticity.)} $W_j(U)$, $l_j$, $\bQ_k(U)$, $Q'_k(U)$ and $D_{W_k(U)}$ are holomorphic in $A'$ on $\{|A'|_\infty<\delta_{AB}\}$. The nonlinear part $\Cav_k(A)$ at the reference $U$ is holomorphic in $(A',A)$ jointly.
\end{enumerate}
\end{proposition}

\begin{proof}\hasdetail{prop:AB}
(i) is Lemmas~\ref{lem:AB-rep} and~\ref{lem:AB-bounds}.

(ii) The proof of Lemma~\ref{lem:one-step-expansion} uses the reality of $W$ in four places; everything else is algebraic, uses Lemma~\ref{lem:avg-transport}(a), which allows complex $W$, or uses only analyticity and Cauchy estimates.
\begin{enumerate}[label=(\arabic*)]
\item \emph{Analyticity.} Relative transports along the contours of at most $2L$ bonds of $\T^{(j-1)}$ used in the step are products of at most $2L$ factors $e^{\ii\delta R(W(\gamma'))B(b)}$, with $\gamma'$ an initial piece. By Lemma~\ref{lem:AB-bounds}(c), $\|R(W(\gamma'))\|\le e^{\kappa_{\rm AB}\delta_j\delta_c}\le2$ instead of $1$. By the same lemma, $W(\Gamma)=e^{\ii Z_\Gamma}W^r_{j-1}(\Gamma)$ with $|Z_\Gamma|\le4C_{\rm AB}\delta_j\delta_c$, so the families $\{W(\Gamma)\}$ lie in $\cD_{C(a\delta_j^2+\delta_j\delta_c)}$, and the families $\{(XW)(\Gamma)\}=\{X_W(\Gamma)W(\Gamma)\}$ lie in $\cD_{C(c_2+a\delta_j^2+\delta_j\delta_c)}\subset\cD_{r_M/2}$. The averaged relative contour variables $X_W(y,x)=\Av(\{X_W(\Gamma)W(\Gamma)W(y,x)^{-1}\})$ are bounded by Lemma~\ref{lem:group-average}(e), as in the real case. The factor $2$ is absorbed by $c_2\to\frac12c_2$. By Lemma~\ref{lem:AB-bounds}(d), the arguments of all logarithms are within $C(c_2+a\delta_j^2+\delta_j\delta_c)$ of $1$, and $\kappa(\tau_X(y))\le2$. The Cauchy estimates for $c_j$, $c^\tau_j$ and $c^p_j$ are then unchanged.
\item \emph{The loops.} In the linear part, $|H_x|,|\Eloop|\le16a\delta_j^2$ is replaced by $16a\delta_j^2+C_{\rm loop}\delta_j\delta_c$, by Lemma~\ref{lem:AB-bounds}(d). The bounds $|\mathfrak f(\operatorname{ad}H_x)-1|\le C|H_x|$, $|\mathfrak e(\operatorname{ad}\Eloop)-1|\le C|\Eloop|$ and $|R(e^{\ii \Eloop})-R(h_x)|\le C(|\Eloop|+|H_x|)$ hold for complex $H_x,\Eloop$ of size at most $1$. Indeed $\mathfrak e$ is entire, $\mathfrak f$ is analytic for $\|\operatorname{ad}\|<2\pi$ (Lemma~\ref{lem:one-step-expansion}), and $\|\operatorname{ad}H\|\le2|H|\le2$. Moreover $W_j(c)=e^{\ii \Eloop}W([c])$, so $\|R(h_xW([c]))-R(W_j(c))\|=\|(R(h_xe^{-\ii \Eloop})-1)R(W_j(c))\|\le C(|\Eloop|+|H_x|)$, since $\|R(W_j(c))\|\le2$.
\item \emph{The tube average and the weights.} In $q_jB(c)$ every value $B(b)$ is multiplied by $R(W(c_-,x))R(W(\gamma'))$, a transport along an averaged contour variable times one along an initial piece $\gamma'$ of $[x,x(c)]$. By Lemma~\ref{lem:AB-bounds}(c) it has norm at most $e^{\kappa_{\rm AB}\delta_j\delta_c}$, and the scalar weights are those at $W=1$. The Cauchy--Schwarz inequality gives the bound on $q_j$. The weights of the averaged relative transports at the complex families satisfy $\sum_\Gamma|\omega_\Gamma|\le1+C_M\max_\Gamma|H_\Gamma|\le1+C(a\delta_j^2+\delta_j\delta_c)$, by Lemma~\ref{lem:group-average}(a),(d) and (1), and the transports in $B^W(\Gamma)$ have norm at most $e^{\kappa_{\rm AB}\delta_j\delta_c}$. So $|\lambda_W(y,x)[B]|\le2(L-1)\big(1+C(a\delta_j^2+\delta_j\delta_c)\big)e^{\kappa_{\rm AB}\delta_j\delta_c}\max|B|$, and $|\gamma_jB|\le de^{\kappa_{\rm AB}\delta_j\delta_c}\max|B|$.
\item \emph{The remainder $r_j$.} The algebra of the proof is unchanged: with all factors $\mathfrak e,\mathfrak f,R(h_x),R(e^{\ii \Eloop})$ equal to $1$ and $R(h_xW([c]))=R(W_j(c))$, only $q_jB$ remains. So $r_jB$ is a sum of terms, each containing one of the differences in (2), multiplied by $\Xi$, $\Xi_x$, $B^W([c])$ or $\lambda_W(c_+,x(c))[B]$, which are bounded by $CLe^{\kappa_{\rm AB}\delta_j\delta_c}\max|B|$ by (3), and by the prefactor $L^{-1}$. Hence $|r_jB|\le C(a\delta_j^2+\delta_j\delta_c)\max|B|$.
\end{enumerate}
Part~(f) follows as in the real case, from the factorization
\[
\overline X(c)=\tau_X(c_-)\,\tilde{\mathfrak a}_j(X)(c)\,R(W_j(c))\big(\tau_X(c_+)^{-1}\big),
\]
in which the linear parts add, and from the Cauchy estimates. Its statements on the coefficient of $B(b_0(c))$ are algebraic: for $x$ on the axis of $c$ the family $G(c_-,x)$ has one element, so $W(c_-,x)=W([c_-,x])$ for every $W$.

(iii) \emph{(a), (b).} By (ii), $|l_jB|_\infty\le\big(1+C(a\delta_j^2+\delta_j\delta_c)\big)|B|_\infty$ and $|c_j(B)|_\infty\le C\delta_j|B|^2_\infty$. The induction of the proof of Proposition~\ref{prop:Q-expansion} gives, by \eqref{eq:AB-scales} and $\sum_i\delta_i^2\le\frac43$,
\[
|B_j|_\infty\le|A|_\infty\prod_{i\le j}\big(1+C(a\delta_i^2+\delta_i\delta_c)+2C\delta_i|A|_\infty\big)\le|A|_\infty e^{C(2a+2\delta_c+4|A|_\infty)}\le2|A|_\infty
\]
for $a,|A|_\infty\le c'_3$ and $\delta_c<\delta_{AB}$. The bounds on $\Cav_j$ and the locality follow as there.

\emph{(c).} We run the proof of Proposition~\ref{prop:Q-expansion}(c) with nonnegative majorants. For a scale-$(m-1)$ field $v$ write $|v|$ for the function $b\mapsto|v(b)|$. In $q_mv(c')$ each value $v(b)$ is multiplied by one transport of norm at most $e^{\kappa_{\rm AB}\delta_m\delta_c}$, by (ii)(3), and the scalar weights are those at $W=1$. So $|q_mv|\le e^{\kappa_{\rm AB}\delta_m\delta_c}q^{(1)}_m|v|$ pointwise, where $q^{(1)}_m$ is the tube average at $W=1$. By (ii)(4), $|r_mv|\le r^{\rm maj}_m|v|$, where $r^{\rm maj}_m$ is a nonnegative kernel with the support of the real case and at most $C_L$ entries per row, each bounded by $C(a\delta_m^2+\delta_m\delta_c)$; it already accounts for the complex transports inside $r_m$. Hence
\[
|l_j\cdots l_{i+1}v|\le e^{\kappa_{\rm AB}\delta_c\sum_m\delta_m}\big(q^{(1)}_j+r^{\rm maj}_j\big)\cdots\big(q^{(1)}_{i+1}+r^{\rm maj}_{i+1}\big)|v| .
\]
The prefactor is at most $e^{2\kappa_{\rm AB}\delta_c}$, by \eqref{eq:AB-scales}. The product of majorants is the real case with unitary transports, except that the coefficients of the insertions are $C(a\delta_m^2+\delta_m\delta_c)$ instead of $Ca\delta_m^2$. The steps ``Transports'' and ``Insertions'' of that proof therefore apply verbatim, with $a\delta_m^2$ replaced by $a\delta_m^2+\delta_m\delta_c$. They show that the product is of type
\[
\prod_m\big(1+C_L(a\delta_m^2+\delta_m\delta_c)\big)\le\prod_m(1+C_La\delta_m^2)\prod_m(1+C_L\delta_m\delta_c)\le(1+Ca)\,e^{C_L\delta_c}.
\]
In the step ``First derivatives'', $dc_m(B)$ plays the role of $r_m$, with coefficients $C\delta_m|A|_\infty$ by the Cauchy estimates of (ii), and the same majorant argument applies. In the step ``Second derivatives'', the type bounds just proved replace those of the real case. So both steps hold with constants multiplied by $e^{C_L\delta_c}$. The consequences follow from the kernel bounds as before.

\emph{(d).} The values of $\mathsf w_k$ at the points of $\T$ are given by Theorem~\ref{thm:correspondence}(c), which is the recursion of Definition~\ref{def:adapted} and holds for complex relative configurations, and by (ii); this gives the same bounds. At fine points, the extensions $\mathcal E^{(i)}$ transport along chains of averaged contour variables, with norm at most $e^{\kappa_{\rm AB}\delta_c}$ instead of $1$, by (iv); the averaged relative contour variables $(X_i)_W(y,x)$ are bounded by Lemma~\ref{lem:group-average}(e), as in the real case.

(iv) For every $Y\in G^c$, $R(Y)^T=R(Y^{-1})$ with respect to $\mathrm b$, since $\operatorname{tr}(YXY^{-1}Z)=\operatorname{tr}(XY^{-1}ZY)$. Hence $q'^{\,T}_jg(x)=R(W_{j-1}(U)(y,x))^{-1}g(y)$ and $q'_jq'^{\,T}_j=I$, as in the proof of Proposition~\ref{prop:linearized-averages}(a). That proof uses only the volume factors, not unitarity. The Hermitian identity $Q'_kQ'^*_k=I$ need not hold for complex $U$; the $\mathrm b$-transpose identity always holds. Locality (d) holds because Lemma~\ref{lem:one-step-basic}(c) holds verbatim for complex configurations, and Lemma~\ref{lem:one-step-expansion}(a) holds by (ii). Covariance (f) follows from Lemma~\ref{lem:one-step-basic}(b) with real $g$ and \eqref{eq:avg-contour-gauge}.

The transport $T_U(y,x)$ along the chain of averaged contour variables is a product over $m<k$ of $R(W_m(U)(x_{m+1},x_m))$, where $x=x_0\in B_1(x_1)$, \dots, $x_{k-1}\in B_k(y)$ is the chain of $x$, and $x_k=y$. By Lemma~\ref{lem:AB-bounds}(c) at index $m+1$, $W_m(U)(x_{m+1},x_m)=e^{\ii Z_m}W_m(U_r)(x_{m+1},x_m)$ with $|Z_m|\le4C_MC_{\rm AB}\delta_{m+1}\delta_c$. By \eqref{eq:Ad-bound}, \eqref{eq:AB-scales} and the unitarity of $R(W_m(U_r)(x_{m+1},x_m))$, $\|T_U\|\le\exp\big(8C_MC_{\rm AB}\delta_c\sum_m\delta_{m+1}\big)\le e^{\kappa_{\rm AB}\delta_c}$. Telescoping the product gives $\|T_U-T_{U_r}\|\le e^{\kappa_{\rm AB}\delta_c}-1$.

For (e), compute in the given gauge, extended arbitrarily to $\Tf$; this is allowed by (f). By locality, $\tilde B_j$ on the relevant bonds depends only on $U_r,A'$ in the cube. By Lemma~\ref{lem:avg-local-gauge}, $W_j(U_r)=e^{\ii\delta_jA_{r,j}}$ with $|A_{r,j}|\le\eps+32a\le C\eps$. By Lemma~\ref{lem:AB-bounds}(a),(b) and Lemma~\ref{lem:log-estimates}(c), $W_j(U)=e^{\ii\delta_jA_{0,j}}$ with $|A_{0,j}|\le C(\eps+\delta_c)$. All transports along contours of at most $2L$ bonds of $\T^{(j-1)}$, and, by Lemma~\ref{lem:group-average}(c), those along the averaged contour variables, are then within $C(\eps+\delta_c)\delta_j$ of $1$. By Lemma~\ref{lem:avg-transport}(e), which allows complex $W$, $|\lambda_W(y,x)[B]-\lambda_1(y,x)[B]|\le C(\eps+\delta_c)L\delta_j\max|B|$. With (ii), $l_j-q_j(1)$ is a combination of at most $C_L$ values with coefficients $C_L(\eps+\delta_c)\delta_j$, using $a\delta_j^2\le\eps\delta_j$. The telescoping of the proof of Proposition~\ref{prop:linearized-averages}(e), with the type bounds of (iii), gives the kernel bound $\sum_jC_L(\eps+\delta_c)\delta_je^{C_L\delta_c}\le C_L(\eps+\delta_c)$. The same argument applies to $Q'_k$.

(v) The bound for $Q'_k$ is the bound on the chain transports in (iv). For the unit-lattice derivative, $W_k(U)(b)=e^{\ii\tilde B_k(b)}W_k(U_r)(b)$ with $|\tilde B_k|\le C_{\rm AB}\delta_c$, and \eqref{eq:Ad-bound} applies. For $\bQ_k$ telescope:
\[
\bQ_k(U)-\bQ_k(U_r)=\sum_{j\le k}l_k(U)\cdots l_{j+1}(U)\,\big(l_j(U)-l_j(U_r)\big)\,l_{j-1}(U_r)\cdots l_1(U_r).
\]
By the formula for $l_j$ in the proof of Lemma~\ref{lem:one-step-expansion}, $l_jB(c)$ is a sum of at most $C_L$ terms. Each is a product of a bounded number of factors of the forms $R(W(\gamma))$, $R(W(y,x))$, $\omega_\Gamma$, $R(e^{\ii \Eloop})$, $R(h_x)$, $\mathfrak e(\operatorname{ad}\Eloop)$ and $\mathfrak f(\operatorname{ad}H_x)$, applied to one value $B(b)$, with $\gamma$ a contour of at most $2L$ bonds. Each factor is bounded by $2$. By Lemma~\ref{lem:AB-bounds}(c),(d), each factor changes by at most $C\delta_j\delta_c$ from $U_r$ to $U$. For the weights $\omega_\Gamma=(\sum_{\Gamma'}\mathfrak f(\operatorname{ad}H_{\Gamma'}))^{-1}\mathfrak f(\operatorname{ad}H_\Gamma)$ this follows from the change of $H_\Gamma=\ell(W(\Gamma)W(y,x)^{-1})$, which is at most $C\delta_j\delta_c$ by Lemma~\ref{lem:AB-bounds}(c) and the argument of Lemma~\ref{lem:AB-bounds}(d). So $l_j(U)-l_j(U_r)$ is a combination of at most $C_L$ values with coefficients $C_L\delta_j\delta_c$. With the type bounds of (iii), at $U$ and at $U_r$, the kernel of the sum is bounded by $\sum_jC_L\delta_j\delta_c(1+Ca)e^{C_L\delta_c}\le C_L\delta_c$. The supports are those of Proposition~\ref{prop:linearized-averages}(d). Transposes follow by summing the kernel bounds.

(vi) By Lemma~\ref{lem:AB-bounds}(b), $W_j(U)$ is holomorphic in $A'$. On the domain of (ii), $\delta_j^{-1}\ell(\tilde{\mathfrak a}_j(e^{\ii\delta_{j-1}B})(c))$ is jointly holomorphic in $B$ and in the finitely many bond variables of $W_{j-1}(U)$ on which it depends. Its derivative at $B=0$, which is $l_j$, is therefore holomorphic in $A'$. The rest follows by composition, uniformly in $k$ by the bounds of (iii)--(v).
\end{proof}

\begin{remark}[Why the naive bound fails]\label{rem:AB-outline}
The linear part in Lemma~\ref{lem:AB-bounds}(a) explains why a per-scale bound fails. The derivative at $A'=0$ of the plain iterated average is $\bQ_j-D_{W^r_j}\boldsymbol\Gamma_j$: the gauge-adapted linearization plus a coarse pure-gauge term. The plain one-step relative average, whose linearization is $p_j=l_j-\delta_jD_{W_j}\gamma_j$ (Lemma~\ref{lem:one-step-expansion}(f)), has a sup-norm loss of up to a factor $2d+1$, which comes entirely from that pure-gauge term. The term does not accumulate, since $\boldsymbol\Gamma_j$ is a sum of $\delta_{i+1}$-weighted contour averages.
\end{remark}

\begin{remark}[Scope of Proposition~\ref{prop:AB}]\label{rem:AB-remains}
\begin{enumerate}[label=(\arabic*)]
\item \emph{Not proved, and not needed below.}
  \begin{itemize}
  \item A curvature bound for $W_j(U)$ of order $a+\delta_c$ in field-strength units for $j<k$, i.e.\ a complex analogue of Lemma~\ref{lem:one-step-regularity} and Proposition~\ref{prop:avg-regular}. We only have $|H_x(U)|\le16a\delta_j^2+C_{\rm loop}\delta_j\delta_c$. This is weaker by $\delta_j^{-1}$ in the $\delta_c$-term, but summable over scales. At the unit scale the two agree: $|F_{W_k(U)}|\le2a+C\delta_c$, from Lemma~\ref{lem:AB-bounds}(c) with $n=4$ and the Schwarz-lemma argument of (d).
  \item Derivative or H\"older bounds of $\tilde B_j$ for $0<j<k$.
  \item Theorem~\ref{thm:correspondence}(a),(b),(e),(f) and Lemma~\ref{lem:axial}(b),(c) at complex configurations. They are stated for real relative configurations; only Theorem~\ref{thm:correspondence}(c),(d) and Lemma~\ref{lem:axial}(a) hold for complex ones.
  \end{itemize}
  Lemma~\ref{lem:one-step-expansion}, Proposition~\ref{prop:Q-expansion} and Proposition~\ref{prop:linearized-averages} use the backgrounds only through transports along contours of bounded length and along averaged contour variables, and through the weights of the averaged relative transports. At complex references Part~X uses the averaging only through the objects of Definition~\ref{def:adapted} and their linearizations, which are defined by the recursion; through Theorem~\ref{thm:correspondence}(c), for the values of $\mathsf w_j$ at the coarse points; and through the exact identities of step~(iii) of the proof of Proposition~\ref{prop:complex-ext}, which are continued from the real points by the identity theorem. The averaged axial gauge and Theorem~\ref{thm:correspondence}(a),(b) enter Part~X only through Part~V, at real references (Theorem~\ref{thm:weak-bg}(c), Lemma~\ref{lem:critical-current}). So nothing is assumed about the third item. If a later argument needs any of the items above, it must be supplied separately.
\item \emph{Structure of the proof.} Proposition~\ref{prop:AB}(ii)--(v) is a line-by-line modification of the proofs in Section~\ref{sec:adapted}, with the type bookkeeping of Proposition~\ref{prop:Q-expansion}(c). The new inputs, Lemmas~\ref{lem:AB-cov}--\ref{lem:AB-bounds}, are self-contained given Proposition~\ref{prop:Q-expansion} at real references.
\item \emph{Downstream.}
  \begin{itemize}
  \item The bounds $|\bQ A(c)|\le C\|A\|_{S_c}$ and $|Q'\lambda(y)|\le C\|\lambda\|_{B(y)}$, and the sup-norm bounds of $\bQ^T,Q'^T$, used in Theorem~\ref{thm:ID} and Part~IX, follow from (iii) and (iv) with constants $e^{C_L\delta_c}\le2$.
  \item The exact identities ($E=0$, Proposition~\ref{prop:linearized-averages}(b), and those listed in step~(iii) of the proof of Proposition~\ref{prop:complex-ext}) are not extended on all of $\{|A'|_\infty<\delta_{AB}\}$: there only the sup norm of $A'$ is controlled, so the real points need not be regular backgrounds. They hold on the ball $\mathcal B$ of that proof, whose real points are regular, by (vi) and the identity theorem.
  \end{itemize}
\end{enumerate}
\end{remark}

\subsection{Extension of the results of Part~IX}\label{sec:complex-ext}

\paragraph{Conventions at complex references.} Proposition~\ref{prop:complex-ext} is read as follows.
\begin{itemize}
\item Each result listed in (a)--(d) of the proposition holds at a complex reference of size $\delta_c<\delta_P$ whose real part is $(M,\eps)$-regular with $\eps\le\eps_{\max}$, with $\eps$ replaced by $\eps+C\delta_c$.
\item The $\Rea$-coercivity constants and the decay rates are fixed and do not depend on $\delta_c$: $\frac c2$ with the $c$ of the real reference, and $\frac18c_L$, $\frac1{16}c_L$ at level~(L5); and the rates that the real-reference proofs give when every input decay bound is taken at a quarter of the minimum of the perturbation rate and the two chosen inverse decay rates in Lemma~\ref{lem:perturbed-inverse}.
\item Every other constant, computed at the real reference with these coercivity constants and rates, changes by at most a factor $1+C_{\rm cx}\delta_c\le2$.
\end{itemize}
The constants involved are chosen in the following order (Remark~\ref{rem:parameters}(4),(5), Remark~\ref{rem:domain-constants}):
\[
L\ \to\ \kappa,\delta\ \to\ \rho\ \to\ M_1,M_{\mathfrak N},M'_1\ \to\ M\ \to\ \delta_*,C_W,c''_L\ \to\ M'\ \to\ C_{\rm cx}\ \to\ \bar\delta\ \to\ \eps_{\max}\ \to\ \delta_P .
\]
The thresholds $M_1$, $M_{\mathfrak N}$, $M'_1$ are computed with doubled constants, so they serve at complex references and are fixed before $M$ (Remark~\ref{rem:complex-in-44}). The radius $\delta_*$ and the constant $C_W=C(1+2\delta_*)/\delta_*$ are fixed in the proof of Proposition~\ref{prop:dirichlet-hessian}(c), depending on $d$, $L$, $G$, $\alpha$ and $M$; the constant $c''_L$ is fixed in the proof of Lemma~\ref{lem:coercive-general}, with $C_{\rm cmp}c''_L\le\delta_*$. The constant $C_{\rm cx}$ depends on $L,M,M',\kappa,\delta,\rho,\alpha$ and $G$, and enters only the cap $\bar\delta\le\min(\delta_{AB},1/C_{\rm cx})$. The threshold $\eps_{\max}$ does not depend on $\bar\delta$ or $\delta_P$, and the conditions on $\delta_P$, among them $\delta_P\le\bar\delta$ and $\delta_P\le c''_L$, are listed at the end of the proof of Proposition~\ref{prop:complex-ext}. The domain constants come after $\delta_P$, and (DC4) makes $\delta_c<\delta_P$ at the references of the step (Remark~\ref{rem:complex-ext-use}).

\begin{proposition}[Extension to complex references]\label{prop:complex-ext}
There is $\delta_P>0$, depending on $L,M,\kappa,\rho,\alpha$, on the loss parameter $\delta$ of Part~IX, and on the cap $\bar\delta$ and the threshold $\eps_{\max}$ of Remark~\ref{rem:domain-constants}, such that the following holds for $\delta_c<\delta_P$, at complex references of size $\delta_c$ whose real part is $(M,\eps)$-regular with $\eps\le\eps_{\max}$. The statements hold with $\eps$ replaced by $\eps+C\delta_c$. The $\Rea$-coercivity constants and the decay rates do not depend on $\delta_c$. The coercivity constants are those of Proposition~\ref{prop:reg-complex}, i.e.\ $\frac c2$ with $c$ that of the real reference (Lemma~\ref{lem:CT-stable}), and at level~(L5) $\frac18c_L$ and $\frac1{16}c_L$ (Lemma~\ref{lem:coercive-general}, Proposition~\ref{prop:local-fluct}). The decay rates are those that the real-reference proofs give when every input decay bound is taken at the fixed rate of the differences, a quarter of the minimum of the perturbation rate and the two chosen inverse decay rates in Lemma~\ref{lem:perturbed-inverse} (step~(vii) of the proof). The other constants, computed at the real reference with these coercivity constants and rates, change by at most a factor $1+C_{\rm cx}\delta_c\le2$, where $C_{\rm cx}$ depends on $L,M,M',\kappa,\delta,\rho,\alpha$ and $G$, but not on $\delta_c$, $\bar\delta$, $\eps_{\max}$ or $\delta_P$, and $C_{\rm cx}\bar\delta\le1$:
\begin{enumerate}[label=(\alph*)]
\item Theorem~\ref{thm:ID}(a)--(c), Corollary~\ref{cor:lp-dirichlet}, and Lemmas~\ref{lem:int-factors} and~\ref{lem:endpoint-factors};
\item Lemma~\ref{lem:trace}, Lemma~\ref{lem:cut-term}, Lemmas~\ref{lem:lift} and~\ref{lem:N-tier}, Proposition~\ref{prop:nl-tiers}, Theorem~\ref{thm:tier-walk} at the levels (L1) and (L3), and Theorem~\ref{thm:HC-pointwise}, with $W=0$ at level~(L3);
\item Proposition~\ref{prop:action}(a)--(c). Here $F^s$ and $J=D^TF^s$ are the holomorphic continuations, with $P_{\mathfrak g}$ replaced by its complex-linear extension $P_{\mathfrak g^c}$ and the summand continued without taking real parts. $|F^s|\le C(\eps+\delta_c)$, while $J$ need not be bounded uniformly in $\eta$;
\item Theorem~\ref{thm:weak-bg}(a), (d), and (b) without the realness statement. The fixed point $A_*(s;B)$ is analytic in $A'$, jointly with $(s,B)$.
\end{enumerate}
Not extended, and not needed at complex references: Theorem~\ref{thm:ID}(d), the identity of Lemma~\ref{lem:div-curvature} (its bound on $W_0$ extends; proof of Proposition~\ref{prop:dirichlet-hessian}(c)), Lemma~\ref{lem:critical-current} and Theorem~\ref{thm:weak-bg}(c). The ``real for real arguments'' statements of Propositions~\ref{prop:Q-expansion} and~\ref{prop:action} and of Lemma~\ref{lem:one-step-expansion}(e),(f) are dropped at complex references as well.
\end{proposition}

\begin{proof}\hasdetail{prop:complex-ext}
We go through the places where reality enters.
\begin{enumerate}[label=(\roman*)]
\item \emph{The $L^2$ theory.}
  \begin{itemize}
  \item Proposition~\ref{prop:reg-complex} gives coercivity (as $\Rea$-coercivity), decay and comparison of domains for the Dirichlet operators of Parts~III and~VIII. It was proved for the averaging of Section~\ref{sec:background}; for the gauge-adapted averaging of Part~IV it needs property~(v) of Remark~\ref{rem:general-averaging}, i.e.\ Proposition~\ref{prop:AB}(iv)--(vi) and step~(iii) below. The averaged operators $\boldsymbol\Gamma^X$, $\bH^X$, $\bC^X$ follow by compositions and perturbed inverses.
  \item Part~IX allows complex perturbations from the outset (Section~\ref{sec:weakenable}, Setting).
  \end{itemize}
\item \emph{The local gauge in the interior estimates.} Instead of a real gauge in which $U$ satisfies (B), use the real gauge of $U_r$. There $U=e^{\ii\eta A_0}$ with complex $A_0$ satisfying (B) with parameter $C(\eps+\delta_c)$.
  \begin{itemize}
  \item Lemmas~\ref{lem:cov-expansion} and~\ref{lem:curvature} allow $\mathfrak g^c$-valued $A$.
  \item The absorption, iteration and $L^p$ arguments use only sup norms and quotients, no maximum principle.
  \item The flat estimates (Theorems~\ref{thm:local-regularity}, \ref{thm:lp-local}, Lemmas~\ref{lem:rescaled}, \ref{lem:lp-kernels}) apply to complex-valued fields componentwise.
  \end{itemize}
\item \emph{Algebraic and analytic identities.} Fix $U_r$ and consider $A'$ in the open ball $\mathcal B=\{\text{size}(A')<\delta_P\}$. The size is a norm on the complex fields $A'$, so $\mathcal B$ is convex. Real $A'$ give real configurations $U=e^{\ii\eta A'}U_r$ that are $(M,\eps+C\delta_P)$-regular, with curvature at most $C(\eps+\delta_P)$; this needs $\eps_{\max}+C\delta_P\le\min(\eps_0,c'_3)$. At these real points the operators below are those of Parts~IV--IX at the real background $U$: Proposition~\ref{prop:AB} defines all objects at the reference $U$ through $W_j(U)=U^{(j)}$ by the formulas of Section~\ref{sec:adapted}, and inverses are unique. So every identity proved in the note for real regular backgrounds of small curvature holds at the real points of $\mathcal B$. These form a totally real subspace of maximal dimension of the convex set $\mathcal B$. So an identity between two holomorphic functions on $\mathcal B$ that holds at the real points holds on all of $\mathcal B$, by the identity theorem. This is used in two stages.
  \begin{itemize}
  \item[(iii-a)] Identities between objects of Proposition~\ref{prop:AB} and of the action. These are holomorphic on $\mathcal B$ by Proposition~\ref{prop:AB}(vi), and since the action is a polynomial in the entries of $U$ and $U^{-1}$. They are the exact covariance $E=0$ (Proposition~\ref{prop:linearized-averages}(b)); $J=D^TF^s$ (Lemma~\ref{lem:current}(a)); and $\ip{\rho^{(1)}(A)}{F^s}=0$ in the proof of Proposition~\ref{prop:action}. After this stage the pair $(Q'_k,\bQ_k)$ has property~(v) of Remark~\ref{rem:general-averaging} on $\mathcal B$, and step~(i) applies. Stage~(iii-a) holds with the same proof on every such ball whose real points satisfy $\eps+C\cdot(\text{radius})\le\min(\eps_0,c'_3)$; in particular on the ball of radius $2\delta_*$ of the proof of Proposition~\ref{prop:dirichlet-hessian}(c).
  \item[(iii-b)] Identities that involve inverses. After step~(i), the operators of Parts~III, VIII and~IX are holomorphic on $\mathcal B$, as compositions and inverses of holomorphic families with uniformly bounded inverses, and uniformly convergent series of such. The identities are those of Lemma~\ref{lem:domain-identities} used in Lemma~\ref{lem:local-system}; $\bG DR=DG'R$ (Lemma~\ref{lem:Api-identities}(a)) and \eqref{eq:C-identity}; the $\mathrm b$-symmetry $(\bG^Z)^T=\bG^Z$ used in Lemma~\ref{lem:cut-term}; and the identity of Proposition~\ref{prop:one-loop-polymer}(a). For the last, $\log\det\Delta_k$ means $\sum_b\ell_k(b)$, with the integral of Proposition~\ref{prop:det-densities}(c). On $\mathcal B$, $\Rea\Delta_k\ge\frac18c_L$ (Lemma~\ref{lem:coercive-general}; it needs $|J_{U_r}|\le\eps$, which holds at every reference where Proposition~\ref{prop:one-loop-polymer} is applied, Section~\ref{sec:step-quantities}, Setting; its conditions $\eps\le\eps_{\rm cc}$ and those of the proof of Proposition~\ref{prop:dirichlet-hessian}(c) hold since $\eps\le\eps_{\max}$, Remark~\ref{rem:domain-constants}), so the spectrum of $\Delta_k$ lies in $\{\Rea z>0\}$. There the integral is the principal logarithm: it is holomorphic on $\mathcal B$, real at the real points, and its exponential is $\det\Delta_k$.
  \end{itemize}
  No connectedness of the set of all complex references is needed. That set need not be connected: if $G$ has a $U(1)$ factor, real regular backgrounds with different fluxes exist.
\item \emph{Unitarity of transports and the averaging.} This is Proposition~\ref{prop:AB}.
\item \emph{The action.} The proof of Proposition~\ref{prop:action} uses a gauge with $|A_0|,|\partial A_0|\le\eps$, $|R(H_p)-1|\le C\eps\eta^2$ and $|F^s|\le C\eps$. These hold with $C(\eps+\delta_c)$, using the holomorphic summand and $P_{\mathfrak g^c}$. The current $J$ is not bounded uniformly in $\eta$. In the weakened background map it enters only as $(D\bH(s))^T\mathsf F$ (Section~\ref{sec:weak-background}).
\item \emph{The trace inequality.} Lemma~\ref{lem:trace} is stated for complex perturbations.
\item \emph{The tiers.} Corollary~\ref{cor:lp-dirichlet} and Lemmas~\ref{lem:int-factors} and~\ref{lem:endpoint-factors}, also with $\varphi\equiv1$, follow from (i), (ii) and (iv) as at real backgrounds. In particular, Lemma~\ref{lem:local-sobolev} and Theorem~\ref{thm:lp-local} are flat estimates, applied componentwise; the normal form holds with complex coefficients of size $C(\eps+\delta_c)$, by (ii); and the bounds $|Q'\lambda(y)|\le C\|\lambda\|_{B(y)}$, the sup bounds of $Q'^T$ and the block-locality of $Q'^TQ'$ hold by (iv). Lemma~\ref{lem:cut-term} follows from (i) and (vi). Theorem~\ref{thm:tier-walk} at level~(L1) has only local steps, which Lemmas~\ref{lem:int-factors} and~\ref{lem:endpoint-factors} bound. Lemmas~\ref{lem:lift} and~\ref{lem:N-tier} and Proposition~\ref{prop:nl-tiers} use, besides these:
  \begin{itemize}
  \item the walks of Part~IX, which allow complex perturbations from the outset, and the comparison of domains, by (i);
  \item Proposition~\ref{prop:scalar-interior}, for the differences of local propagators, by (ii);
  \item the bounds on $Q'$, $Q'^T$ and $Q'^TQ'$ just listed, by (iv);
  \item Lemmas~\ref{lem:trace} and~\ref{lem:cut-term}, for an external $D$, by (vi);
  \item Lemma~\ref{lem:curvature} with $\mathfrak g^c$-valued $A$, which gives $|c_{\mu\nu}|+[c_{\mu\nu}]_\beta\le C(\eps+\delta_c)$ for the coefficients of $\mathfrak F$ in Lemma~\ref{lem:N-tier}(N2).
  \end{itemize}
  Their proofs apply verbatim, and no new condition on $\delta_c$ arises; the thresholds used are those listed below. \emph{The fixed rates.} For the operators of Parts~II, III and~VIII, Lemma~\ref{lem:CT-stable} gives the hypotheses of Lemma~\ref{lem:CT-form} with the constants $c/2$, $C_1+1$ and $\min(\alpha_1,\gamma/2,1)$, where $c,C_1,\alpha_1$ are those at the real reference and $\gamma$ is the rate of the perturbation. So these operators decay at the rate $\gamma_{\rm c}=\min(\alpha_1,\gamma/2,\frac13,c/(4(C_1+1)))$. By Lemma~\ref{lem:perturbed-inverse}, their differences from the operators at $U_r$ have decay $(C\delta,\gamma_{\rm c}/4)$, with $\delta\le C_{\rm cmp}\delta_c$. At level~(L5) the lower bounds are $\frac18c_L$ and $\frac1{16}c_L$, and the differences again decay at a quarter of the minimum of the perturbation rate and the two chosen inverse decay rates. Every input decay bound below, at real and at complex references, is taken at these quarter rates. At the rate $\gamma_{\rm c}$ itself the constants at a complex reference are only bounded by about twice the real ones.

  \emph{The other constants.} In outline: each constant is a nonnegative combination of products of monotone, positively homogeneous functions of finitely many inputs; at a complex reference the inputs change by at most a factor $1+C_{\rm in}\delta_c$, and convexity in the scaling factor turns this into a change of every constant by at most $1+C_{\rm cx}\delta_c$, level by level. In detail, order the levels of Sections~\ref{sec:nested-walks} and~\ref{sec:tiers}, and the constants of Section~\ref{sec:weak-background}, as in Section~\ref{sec:nested-walks}: $\ell=1,\dots,N$. Each multiplicative constant $K$ of level $\ell$ is $F_K(x_\ell)$, where $x_\ell$ collects the primitive inputs and the constants of the earlier levels, and $F_K$ is a nonnegative combination of products of nonnegative, nondecreasing, positively $1$-homogeneous functions of $x_\ell$. These include sums, products, maxima, Neumann and walk series with nonnegative coefficients, and integrals of such, like the resolvent integral of Corollary~\ref{cor:weak-covariance}. So $t\mapsto F_K((1+t)x_\ell)$ is convex. Put $\sigma_{\rm cx}=2^{1/3}-1$, so that $(1+\sigma_{\rm cx})^3\le2$. For $x'\le(1+t)x$ and $0\le t\le\sigma_{\rm cx}$, convexity gives
  \[
  F_K(x')\le F_K(x)\Big(1+\frac t{\sigma_{\rm cx}}\Big(\frac{F_K((1+\sigma_{\rm cx})x)}{F_K(x)}-1\Big)\Big).
  \]
  At a complex reference the primitive inputs change by at most a factor $1+C_{\rm in}\delta_c$: the averaging kernels by $e^{C_L\delta_c}\le1+2C_L\delta_c$ (Proposition~\ref{prop:AB}(iii)--(v)); the operators of Parts~II, III and~VIII by the differences above; and $W_2^X$ by $C_{\rm cmp}C_W\delta_c$ (proof of Proposition~\ref{prop:dirichlet-hessian}(c)). Each input bound $K$ is taken free of $\eps$ (with $\eps\le1$; e.g.\ $\|W_2^X\|\le C$), so an additive change $C\delta_c$ is a factor $1+(C/K)\delta_c$; for inverses, Lemma~\ref{lem:perturbed-inverse} gives the change relative to the bound, also at level~(L5), where $C_G<1$. Put $C^{(\ell)}=(C^{(\ell)}_{\rm in}/\sigma_{\rm cx})\max\big(1,\max_K(F_K((1+\sigma_{\rm cx})x_\ell)/F_K(x_\ell)-1)\big)$, the maximum over the constants of level $\ell$, with $C^{(\ell)}_{\rm in}=\max(C_{\rm in},C^{(1)},\dots,C^{(\ell-1)})$. Let $C_{\rm cx}=\max_\ell C^{(\ell)}$. For $\delta_c\le1/C_{\rm cx}$ one has $C^{(\ell)}_{\rm in}\delta_c\le\sigma_{\rm cx}$ for every $\ell$. So, by induction over $\ell$, every input of every level changes by at most a factor $1+\sigma_{\rm cx}$, and every constant by at most $1+C_{\rm cx}\delta_c\le2$. The values $F_K((1+\sigma_{\rm cx})x_\ell)$ are finite. Indeed, the step constants $\theta$ of Lemma~\ref{lem:step-small} and $\theta'$ of Theorem~\ref{thm:tier-walk} are sums of products of at most three inputs of their level, so their values at $(1+\sigma_{\rm cx})x_\ell$ are at most twice those at $x_\ell$. Moreover $M_1$ is computed with the inputs doubled, and $M_1'$ with $\theta'$ replaced by $2\theta'$. So for $M\ge M_1,M_1'$ the walks converge at $(1+\sigma_{\rm cx})x_\ell$, with ratio at most $\frac12$. For the same reason $2\theta'$ bounds $\theta'$ at complex references, and the inputs of $M_1$ at most double there. $C_{\rm cx}$ is built from real-reference constants that hold for all $\eps$ below thresholds fixed before $\bar\delta$, such as $\eps_0$ and $c_5$. So it depends on $L$, $M$, $M'$, $\kappa$, $\delta$, $\rho$, $\alpha$ and $G$, but not on $\bar\delta$, $\eps_{\max}$ or $\delta_P$. The thresholds $M_{\mathfrak N}$, $M_1$ and $M_1'$, fixed in this way before $M$, serve at every complex reference with $\delta_c<\delta_P$. The proofs of Theorems~\ref{thm:tier-walk} and~\ref{thm:HC-pointwise} then apply verbatim.
\item \emph{Analyticity in $A'$.} Every input is holomorphic in $A'\in\mathcal B$, by (i)--(v), and the contraction of Theorem~\ref{thm:weak-bg} is uniform in $A'$. The factorization of Theorem~\ref{thm:weak-bg}(d), including its statement on the dependence on $s$, follows from the same proof: Lemma~\ref{lem:weak-basic}(b) allows complex backgrounds, the averaging operators are local at complex references (Proposition~\ref{prop:AB}(iii),(iv)), and the nonlinear part is local by (c).
\item \emph{Not extended.}
  \begin{itemize}
  \item The identity of Lemma~\ref{lem:div-curvature} differentiates real gauge transformations and needs $|J_U|\le\eps$. It enters only Theorem~\ref{thm:ID}(d), which is not used after Section~\ref{sec:weak-background}.
  \item Lemma~\ref{lem:critical-current} and Theorem~\ref{thm:weak-bg}(c) concern real critical backgrounds, where they are used.
  \end{itemize}
\end{enumerate}
\emph{Smallness.} $\delta_P$ must satisfy the following:
\begin{itemize}
\item $\delta_P\le c''_L$ of Lemma~\ref{lem:coercive-general}, for its coercivity, used in Section~\ref{sec:step-quantities}. Since $C_{\rm cmp}c''_L\le\delta_*$, with $C_{\rm cmp}=4$ (first item of Section~\ref{sec:complex-references}) and $\delta_*$ of the proof of Proposition~\ref{prop:dirichlet-hessian}(c), the perturbations of size $<C_{\rm cmp}\delta_P$, and the ball of twice that size used there for the Cauchy estimate, lie within the radii fixed with $\delta_*$. These are the radii of Proposition~\ref{prop:reg-complex} for $W=0$ and for $W_2$, $W_2^X$, of Proposition~\ref{prop:AB}, of Proposition~\ref{prop:action}, of Proposition~\ref{prop:dirichlet-averaged}(e), and of $E=0$ (step~(iii)). With $C_{\rm AB}c''_L\le c'_L$ the condition gives $C_{\rm AB}\delta_P\le c'_L$; moreover $C_{\rm AB}\delta_P\le cL^{-d-1}$, for the elimination of Lemma~\ref{lem:eliminate} at a complex base (Corollary~\ref{cor:covariance}(c), with $|\mathfrak r_c|\le\frac12$);
\item $\eps_{\max}+C\delta_P\le\min(\eps_0,c'_3,\eps_{\rm ID},\eps_{\rm cc},\eps_S,\eps_L,\eps'_S,\eps'_L,\eps',\frac1{20})$, with $\eps_S$, $\eps_L$ the thresholds of Propositions~\ref{prop:scalar-interior} and~\ref{prop:system-interior}, $\eps'_S$, $\eps'_L$ those of Propositions~\ref{prop:scalar-lp} and~\ref{prop:system-lp}, and $\eps'$ that of Corollary~\ref{cor:lp-dirichlet}; moreover $\eps_{\max}+C\delta_P$ satisfies the smallness condition on $\eps$ depending on $M$ of Section~\ref{sec:nested-walks}, and $C(\eps_{\max}+\delta_P)\le1$, for the normal form of Lemma~\ref{lem:cov-expansion}, used in Lemmas~\ref{lem:endpoint-factors} and~\ref{lem:trace};
\item $C(\eps_{\max}+\delta_P)\le c_5$ of Proposition~\ref{prop:action};
\item $\eps_{\max}+C\delta_P\le\eps_P$, with $\eps_P$ of Theorem~\ref{thm:weak-bg} computed with the multiplicative constants doubled;
\item at complex references the condition $(1+K_H)r_P\le c_3$ in the proof of Lemma~\ref{lem:Ds} becomes $(1+2K_H)r_P\le c'_3$, with $r_P$ computed with the multiplicative constants doubled;
\item $\delta_P\le c_{\rm DC}r_ML^{-1}$, the condition on $\delta_P$ in (DC5) of Remark~\ref{rem:domain-constants}. At a complex reference the base of the step is $W=e^{\ii\beta}W_r$ bondwise, with $|\beta|_\infty\le C_{\rm AB}\delta_c$ (Proposition~\ref{prop:AB}(i)). For $c_{\rm DC}$ small, (DC5) gives $C_{\rm AB}\delta_P\le\frac1{32}c_0L^{-1}$, with $c_0$ of Lemma~\ref{lem:exp-gauge}(c), and $LC_{\rm AB}\delta_P\le c$, with $c$ of Lemma~\ref{lem:avg-transport}(f). So the averaged contour variables of the step and the averaged relative transports $\mathfrak G_W$ are defined and holomorphic at the complex base (Remark~\ref{rem:domain-constants}); the analogous condition on the real part is the first condition of (DC5). This is used in Section~\ref{sec:step-quantities}.
\item $\delta_P\le\bar\delta$, so that $C_{\rm cx}\delta_P\le1$ (step~(vii)).
\end{itemize}
At complex references the constant $K_H$ is at most $2K_H$, by step~(vii). Later sections use $2K_H$ in place of $K_H$, and the radii $R_P$, $r_P$ and the constants $C_P$, $\eps_P$, $\varrho_P$ of Theorem~\ref{thm:weak-bg} computed with the multiplicative constants doubled. These bound, respectively are valid for, the corresponding quantities at every complex reference with $\delta_c<\delta_P$, and they do not depend on $\delta_P$. Here $\eps_{\max}$ and the cap $\bar\delta\ge\delta_P$ are those of Remark~\ref{rem:domain-constants}; $\eps_{\max}$ is fixed before $\delta_P$, by the rule stated there, and the proposition holds at every complex reference whose real part is $(M,\eps)$-regular with $\eps\le\eps_{\max}$. Every joint condition above has the form $\eps+C\delta_c\le T$, with $T$ fixed before both constants. The thresholds $M_{\mathfrak N}$ and $M_1'$ of Section~\ref{sec:tiers} are lower bounds for $M$ (Remark~\ref{rem:parameters}(4)). So they are fixed before $M$, and hence before $\bar\delta$, with their constants doubled as in step~(vii).
All of these are conditions on $\delta_P$ after $L,M,\kappa,\rho$, $\bar\delta$ and $\eps_{\max}$ are fixed.
\end{proof}

\begin{remark}[Use in the renormalization group]\label{rem:complex-ext-use}
The references of the step are the points of the graded domain of step $k+1$, read in the units of step $k$ and extended as above. Their complex parts have size $\delta_c\le2C_{\rm ext}(L^{-1}\delta_0+L^{-2}\delta_1)$. So $\delta_c<\delta_P$ is the condition (DC4) on the domain constants, which are chosen after $\delta_P$ (Remark~\ref{rem:domain-constants}); it does not involve the coupling. The domains are those of fixed size of Remark~\ref{rem:domain-constants}, and (DC5) makes the averaged contour variables of the step analytic at these references. By Propositions~\ref{prop:AB} and~\ref{prop:complex-ext}, the new activities produced by the step are analytic functions of the complex background, as the norm requires.
\end{remark}

\begin{remark}[Complex references in Section~\ref{sec:tiers}]\label{rem:complex-in-44}
Proposition~\ref{prop:nl-tiers} and Theorem~\ref{thm:tier-walk} concern real backgrounds. The paragraphs of their statements on complex references fix how the thresholds $M_{\mathfrak N}$ and $M'_1$ are computed, so that the same numbers serve at the complex references of this section. Remark~\ref{rem:domain-constants}, Remark~\ref{rem:parameters}(4) and step~(vii) of the proof of Proposition~\ref{prop:complex-ext} use the same rule. It has four parts.
\begin{itemize}
\item[(a)] \emph{Fixed coercivity constants and rates.} The coercivity constants are the fixed ones of Proposition~\ref{prop:complex-ext}, and every decay bound that enters is taken at the fixed rate of the differences there, a quarter of the minimum of the perturbation rate and the two chosen inverse decay rates in Lemma~\ref{lem:perturbed-inverse}. Neither depends on $\delta_c$, so the same numbers serve at real and at complex references.
\item[(b)] \emph{Doubled constants.} The threshold $M_1$ of Theorem~\ref{thm:weak-inverse}, wherever it enters (Remark~\ref{rem:parameters}(4), Proposition~\ref{prop:nl-tiers}), is computed with the multiplicative constants of Part~IX multiplied by~$2$, and $M'_1$ with $\theta'$ replaced by $2\theta'$. Likewise $2K_H$ replaces $K_H$, and the radii $R_P$, $r_P$ and the constants $C_P$, $\eps_P$, $\varrho_P$ of Theorem~\ref{thm:weak-bg}, and the regularity thresholds that define $\eps_{\max}$, are computed with the multiplicative constants doubled (end of the proof of Proposition~\ref{prop:complex-ext}, Remark~\ref{rem:domain-constants}).
\item[(c)] \emph{Why doubling suffices.} At complex references of size $\delta_c<\delta_P\le\bar\delta\le1/C_{\rm cx}$, every multiplicative constant changes by at most a factor $1+C_{\rm cx}\delta_c\le2$, and the inputs of $\theta'$ by at most a factor $1+\sigma_{\rm cx}$, where $\sigma_{\rm cx}=2^{1/3}-1$, so that $(1+\sigma_{\rm cx})^3\le2$. Since $\theta'$ is a sum of products of at most three inputs, $2\theta'$ bounds $\theta'$ there (Proposition~\ref{prop:complex-ext}, step~(vii) of the proof).
\item[(d)] \emph{The order.} The numbers of (a) and (b) do not depend on $\bar\delta$, $\eps_{\max}$ or $\delta_P$. So $M_{\mathfrak N}$, $M_1$ and $M'_1$ are lower bounds for $M$, fixed before $M$ (Remark~\ref{rem:parameters}(4)). The constant $C_{\rm cx}$ may depend on $M$ and $M'$, but it enters only the cap $\bar\delta\le\min(\delta_{AB},1/C_{\rm cx})$, which is chosen after them; then come $\eps_{\max}$ and $\delta_P\le\bar\delta$ (Remark~\ref{rem:domain-constants}).
\end{itemize}
\end{remark}
